\section{Conclusion}

Transcribing the WireGuard handshake once and giving it to two provers
produced three findings a single-tool analysis would not have, and one
property of the artifact itself.

\para{The pre-shared key's contribution became an experiment.} Two models
differing in one line, one proved and one refuted against an attacker holding
a discrete-logarithm oracle. What was previously an assertion about
\texttt{IKpsk2} is now a control with a recorded verdict on each arm
(Table~\ref{tab:psk}).

\para{KCI resistance turned out to have a location.} It is not a property of
the handshake as a whole but of the interval between committing transport keys
and accepting traffic, and the static-ephemeral term \(se\) is what holds that
interval open. Finding this required asking the question at two levels, and
answering it required the prover that ProVerif's saturation had defeated.

\para{Post-quantum cost is not one number.} Following KEM parameter sizes
through to path MTU shows an in-band ML-KEM-768 handshake exceeding the IPv6
minimum by \wgOvershoot\,B. The cost appears as packet size, path
reachability, and assumption strength simultaneously, and a design that
addresses one of the three has not addressed the problem.

\para{The disagreements are the informative part.} The \wgNonterm{}
properties nobody could decide sit in Table~\ref{tab:results} alongside the
rest. A
result matrix in which every cell is green says less about a pair of provers
than one in which the disagreements are visible. That this remains true of
every future run is not a promise but a build step: the tables, figures and
quoted numbers in this paper are regenerated from the provers' own output
each time the document is compiled.
